\chapter{Introduzione}

In questo primo capitolo viene spiegata l'organizzazione del progetto.

%%%%%%%%%%%%%%%%%%%%%%%%%%%%%%%%%%%%%%%%%%%%%%%%%%%%
\section{Main}
Il main del progetto è rappresentato dal Jupyter notebook \texttt{main.ipynb}. È stata presa questa decisione per aumentare la leggibilità delle funzioni chiamate per arrivare ai risultati. La struttura del notebook infatti è la stessa di questa relazione.

Questo documento è pensato per non far eseguire il notebook a chi è interessato ai risultati. Ogni grafico stampato dal codice, derivato dai log su W\&B \cite{wandb} (framework usato durante tutto il progetto per tracciare gli esperimenti) o risultato dei vari esperimenti è riassunto in questa relazione. Con questa idea i vari file \texttt{.py} sono utili non per capire cosa è stato fatto, ma come è stato fatto.

%%%%%%%%%%%%%%%%%%%%%%%%%%%%%%%%%%%%%%%%%%%%%%%%%%%%
\section{Moduli Python}
Le funzioni chiamate nel \texttt{main.py} sono importate da quattro moduli:
\begin{itemize}
    \item \texttt{myModel.py} \\
        Contiene i modelli usati nel progetto e sono classi che estendono \texttt{nn.Module}. Spesso si utilizzano gli stessi blocchi di base, definiti in \texttt{modelBlock.py}.
    \item \texttt{coreFunction.py} \\
        Contiene le funzioni responsabili dell'addestramento, dell'estrazione dell'accuracy, della pipeline usata durante tutto il progetto e del feature extractor.
    \item \texttt{utilityFunction.py} \\
        Contiene funzioni ausiliarie per rendere il codice più leggibile e riutilizzabile. Sono funzioni chiamate più volte e che svolgono compiti precisi.
\end{itemize}